\documentclass[10pt,twocolumn]{article}

\usepackage[margin=1in]{geometry}
\usepackage{amsmath,amssymb}
\usepackage{algorithm}
\usepackage{algpseudocode}
\usepackage{booktabs}
\usepackage{graphicx}
\usepackage{xcolor}
\usepackage{enumitem}
\usepackage{parskip}
\usepackage{titlesec}
\usepackage{fancyhdr}
\usepackage{mdframed}
\usepackage{array}
\usepackage{multirow}
\usepackage{hyperref}

\hypersetup{colorlinks=true, linkcolor=blue, urlcolor=blue}

\titleformat{\section}{\large\bfseries}{}{0em}{}[\titlerule]
\titleformat{\subsection}{\normalsize\bfseries}{}{0em}{}

\pagestyle{fancy}
\fancyhf{}
\rhead{Heuristic v5 --- Utility-Density Refine Mapping}
\lhead{USRT Scheduling}
\cfoot{\thepage}

\definecolor{boxblue}{RGB}{230,240,255}
\definecolor{boxgray}{RGB}{245,245,245}
\definecolor{highlight}{RGB}{255,245,210}

\newmdenv[backgroundcolor=boxblue, linewidth=1pt, linecolor=blue!40,
          innerleftmargin=6pt, innerrightmargin=6pt,
          innertopmargin=4pt, innerbottommargin=4pt]{bluebox}

\newmdenv[backgroundcolor=highlight, linewidth=1pt, linecolor=orange!60,
          innerleftmargin=6pt, innerrightmargin=6pt,
          innertopmargin=4pt, innerbottommargin=4pt]{yellowbox}

\newmdenv[backgroundcolor=boxgray, linewidth=0.5pt, linecolor=gray!60,
          innerleftmargin=6pt, innerrightmargin=6pt,
          innertopmargin=4pt, innerbottommargin=4pt]{graybox}

\algnewcommand\algorithmicforeach{\textbf{for each}}
\algdef{S}[FOR]{ForEach}[1]{\algorithmicforeach\ #1\ \algorithmicdo}

\begin{document}

% ─────────────────────────────────────────────────────────────────────────────
\twocolumn[{%
\begin{center}
  {\LARGE\bfseries Heuristic v5: Utility-Density Aware Mapping Refinement}\\[0.4em]
  {\large\itshape Post-SPS Balancing of Optional-Segment Utility Potential}\\[0.4em]
  {\normalsize Technical Report --- USRT Scheduling System}
\end{center}
\vspace{0.5em}
\hrule
\vspace{1em}
}]

% ─────────────────────────────────────────────────────────────────────────────
\section{Motivation and Problem}

\subsection{The USRT Problem}
The Utility-Aware Scheduling with Real-Time constraints (USRT) problem maps
periodic real-time tasks to processors. Each task $T_i$ has:
\begin{itemize}[nosep]
  \item A \textbf{mandatory segment} with execution time $e_{i,0}$, which
        \emph{must} always complete before the deadline.
  \item $N_i^{seg}$ \textbf{optional segments} with incremental times
        $e_{i,1},\ldots,e_{i,N_i^{seg}}$. Executing optional segments yields
        utility, but is not required.
  \item A utility rate $u_i$ (reward per unit of optional execution time).
\end{itemize}
The objective is to maximise total utility subject to timing feasibility (EDF)
and a global energy budget $B$.

\subsection{What SPS Does --- and Misses}

SPS (Sum Partial Solutions) is used in Phase~1 to map jobs to processors. Its
load metric is mandatory \textbf{utilisation}:
\[
  \ell_i \;=\; \frac{e_{i,0}}{p_i}
\]
SPS minimises the gap between the most- and least-loaded processor, producing
a mapping where every processor carries approximately equal mandatory load.
This is correct for timing feasibility: by Liu \& Layland, EDF schedules a
processor successfully if and only if $\sum \ell_i \leq 1$.

\begin{yellowbox}
\textbf{The SPS inadequacy.}
SPS knows nothing about optional segments or utility. It can map all
high-utility tasks onto one processor. Even if the energy budget allows their
optional segments to run, the \emph{timing constraint} on that congested
processor prevents it. Optional segment time slots compete, and many are lost.
The other processor has free time but nothing valuable to fill it with.
\end{yellowbox}

\subsection{Heuristic v5 --- the fix}

Heuristic v5 inserts a new \textbf{Phase~1b} (Refine Mapping) immediately
after Phase~1 (SPS). It adjusts \emph{which jobs sit on which processor}
to spread utility potential evenly, without breaking the feasible mandatory
mapping SPS produced.

Two algorithms are provided:
\begin{itemize}[nosep]
  \item \textbf{v5a} --- Step 2a: single linear scan over all doublets.
  \item \textbf{v5b} --- Step 2b: iterative one-swap-one-scan, always
        targeting the most overloaded processor.
\end{itemize}
Phases 2--6 (left-shift, energy, freq-scaling, greedy optional segments,
pairwise swap) are identical to Heuristic v4.

% ─────────────────────────────────────────────────────────────────────────────
\section{Key Definitions}

\subsection{Cumulative Execution Time}
Let
\[
  c_{i,k} \;=\; \sum_{q=0}^{k} e_{i,q}, \quad k \in [0, N_i^{seg}]
\]
so $c_{i,0} = e_{i,0}$ (mandatory only) and $c_{i,N_i^{seg}}$ is the total
execution time if all segments run (at $f_{\max}$).

\subsection{Utility Density of a Task}
\begin{bluebox}
\[
  ud(i) \;=\; \frac{u_i \;\times\; \bigl(c_{i,N_i^{seg}} - c_{i,0}\bigr)}{p_i}
\]
\end{bluebox}

\noindent\textit{Plain English}: if this task could always run all its optional
segments, how much utility would it generate per unit of clock time?

The numerator $u_i \times (c_{i,N_i^{seg}} - c_{i,0})$ is the maximum utility
gain from a single job. Dividing by $p_i$ converts it to a rate (utility per
unit time). Every job of task $T_i$ has the same $ud(i)$.

\subsection{Processor Utility Density}
\[
  UD(x) \;=\; \sum_{(i,j)\;\text{mapped to}\;x} ud(i)
\]
This is the total utility potential sitting on processor $x$. The refinement
goal is to make $UD(x)$ approximately equal across all processors.

\subsection{Mandatory Utilisation}
\[
  \mu(i) \;=\; \frac{e_{i,0}}{p_i}, \qquad
  MU(x) \;=\; \sum_{(i,j)\;\text{mapped to}\;x} \mu(i)
\]
SPS balances $MU(x)$. The refinement must not disturb this balance.

% ─────────────────────────────────────────────────────────────────────────────
\section{Doublets}

A \textbf{doublet} is any pair of jobs $(a, b)$ drawn from the full job set
(one job from any task, one from any other task, or the same task).
All $\binom{N_{\text{jobs}}}{2}$ pairs are formed and sorted by
\textbf{ascending} $|ud(a) - ud(b)|$.

\begin{graybox}
\textit{Why this order?} Jobs with similar $ud$ are the most interchangeable
--- swapping them between processors barely changes the UD distribution.
Jobs with very different $ud$ are the ``big moves''. Sorted ascending,
the list moves from safe, small adjustments up to large disruptive ones.
Both 2a and 2b scan this list from the front.
\end{graybox}

Doublets are built once before any swaps and reused throughout the algorithm.
Each pair is considered at most a bounded number of times.

% ─────────────────────────────────────────────────────────────────────────────
\section{Swap Acceptance Conditions}

Any proposed swap of job $a$ (currently on processor $P_a$) with job $b$
(currently on processor $P_b$) must pass \emph{all three} of the following
conditions before being committed.

\subsection*{Condition~1 --- UD imbalance decreases}

Compute hypothetical processor UD values after the swap:
\[
  UD'(P_a) = UD(P_a) - ud(a) + ud(b)
\]
\[
  UD'(P_b) = UD(P_b) - ud(b) + ud(a)
\]
\textbf{For Step 2a (pair-level check):}
\[
  |UD'(P_a) - UD'(P_b)| \;<\; |UD(P_a) - UD(P_b)| - \varepsilon
\]
The imbalance between these two specific processors must strictly decrease.

\textbf{For Step 2b (global max-UD check):}
\[
  \max\bigl(UD'(P_1),\; UD'(P_b)\bigr) \;<\; UD(P_1) - \varepsilon
\]
where $P_1$ is the globally highest-UD processor. The new maximum UD across
both affected processors must be strictly less than the current maximum.
This also prevents cycling (each accepted swap strictly reduces the global max).

\subsection*{Condition~2 --- Mandatory utilisation balance does not worsen}

\begin{yellowbox}
\textbf{Why this guard is needed.}
If a swap moves a high-$e_m$ job onto an already-loaded processor,
mandatory utilisation becomes unbalanced. The overloaded processor then
has little timing slack left for optional segments, even though UD is
``balanced'' on paper. Phase~5 finds fewer optional segments can fit,
and total utility drops. This was the root cause of the ``low utility''
failure before the guard was added.
\end{yellowbox}

Compute the mandatory utilisation imbalance before and after:
\[
  \delta_{\text{cur}} = |MU(P_a) - MU(P_b)|
\]
\[
  \delta_{\text{new}} = \bigl|(MU(P_a) - \mu(a) + \mu(b)) -
                               (MU(P_b) - \mu(b) + \mu(a))\bigr|
\]
Accept only if $\delta_{\text{new}} \leq \delta_{\text{cur}} + \varepsilon$.

\subsection*{Condition~3 --- DBF timing feasibility}

Build hypothetical job lists for both processors after the swap and run
a full mandatory-only DBF check at $f_{\max}$:
\[
  \text{For every window } [t_1, t_2] \text{ on processor } x\text{:}\quad
  \sum_{\substack{(i,j)\;\text{on}\;x\\r_{i,j}\geq t_1,\;d_{i,j}\leq t_2}}
  \!\!e_{i,0} \;\leq\; t_2 - t_1
\]
If either affected processor violates this, the swap is rejected.
This is the hard correctness guarantee: no committed swap ever makes
the mandatory schedule infeasible.

This check is performed on \emph{temporary} copies of the job lists.
The real data structures are not modified. Only if all three conditions
pass does the swap get committed via \textsc{DoSwap}.

% ─────────────────────────────────────────────────────────────────────────────
\section{Algorithm Step 2a --- Linear Doublet Scan (v5a)}

\begin{algorithm}[H]
\caption{Refine Mapping --- Step 2a}
\begin{algorithmic}[1]
\State Compute $ud(i)$, $\mu(i)$ for all tasks
\State Initialise $UD(x)$, $MU(x)$ for all processors from current mapping
\State Build $\mathit{doublets}$ $\leftarrow$ all job pairs sorted by
       ascending $|ud(a)-ud(b)|$
\ForEach{$(a,\,b) \in \mathit{doublets}$}
  \State $P_a \leftarrow \text{mapping}[a]$;\;
         $P_b \leftarrow \text{mapping}[b]$
  \If{$P_a = P_b$} \textbf{continue} \EndIf
  \State Compute $|UD'(P_a)-UD'(P_b)|$ \Comment{Condition 1}
  \If{not improved} \textbf{continue} \EndIf
  \State Compute $\delta_{\text{new}}$ \Comment{Condition 2}
  \If{$\delta_{\text{new}} > \delta_{\text{cur}} + \varepsilon$}
    \textbf{continue}
  \EndIf
  \If{\textsc{SwapFeasible}$(a, b, P_a, P_b)$} \Comment{Condition 3}
    \State \textsc{DoSwap}$(a, b, P_a, P_b)$
    \State Update $UD(P_a)$, $UD(P_b)$, $MU(P_a)$, $MU(P_b)$
  \EndIf
\EndFor \Comment{No break --- every pair is considered once}
\end{algorithmic}
\end{algorithm}

\textbf{Key property of 2a:} the inner loop \emph{never breaks early}. Every
doublet pair is visited exactly once regardless of how many swaps have already
been made. When a swap is committed, $UD$ and $MU$ are updated immediately,
so subsequent pairs see the current state. This is a \emph{distributed}
approach: it fixes many local imbalances in a single left-to-right pass.

% ─────────────────────────────────────────────────────────────────────────────
\section{Algorithm Step 2b --- One-Swap-One-Scan (v5b)}

\begin{algorithm}[H]
\caption{Refine Mapping --- Step 2b}
\begin{algorithmic}[1]
\State Compute $ud(i)$, $\mu(i)$; initialise $UD(x)$, $MU(x)$
\State Build $\mathit{doublets}$ as in Step 2a
\Loop
  \State $P_1 \leftarrow \arg\max_x\; UD(x)$
  \State $P_2 \leftarrow \arg\min_x\; UD(x)$
  \If{$UD(P_1) - UD(P_2) < \varepsilon$} \textbf{break} \EndIf
  \State $\mathit{swapped} \leftarrow \text{false}$
  \ForEach{$(a,\,b) \in \mathit{doublets}$}
    \State $\mathit{hi} \leftarrow$ job with higher $ud$ in the pair
    \State $\mathit{lo} \leftarrow$ the other job
    \If{$\text{mapping}[\mathit{hi}] \neq P_1$} \textbf{continue} \EndIf
    \State $P_b \leftarrow \text{mapping}[\mathit{lo}]$
    \If{$P_b = P_1$} \textbf{continue} \EndIf
    \If{$\max(UD'(P_1),\,UD'(P_b)) \geq UD(P_1)-\varepsilon$}
      \textbf{continue} \Comment{Condition 1}
    \EndIf
    \If{$\delta_{\text{new}} > \delta_{\text{cur}} + \varepsilon$}
      \textbf{continue} \Comment{Condition 2}
    \EndIf
    \If{\textsc{SwapFeasible}$(\mathit{hi},\mathit{lo},P_1,P_b)$}
      \Comment{Condition 3}
      \State \textsc{DoSwap}$(\mathit{hi},\mathit{lo},P_1,P_b)$
      \State Update $UD$, $MU$ for $P_1$ and $P_b$
      \State $\mathit{swapped} \leftarrow \text{true}$
      \State \textbf{break} \Comment{Exit inner loop, re-identify $P_1$}
    \EndIf
  \EndFor
  \If{not $\mathit{swapped}$} \textbf{break} \EndIf
\EndLoop
\end{algorithmic}
\end{algorithm}

\textbf{Key property of 2b:} the inner loop breaks as soon as the
\emph{first} valid swap is found. Control returns to the outer loop, which
re-identifies $P_1$ from the updated $UD$ values. The next inner scan
restarts from the beginning of the doublet list with this new target.

\textbf{Termination guarantee:} Condition~1 for 2b requires
$\max(UD'(P_1), UD'(P_b)) < UD(P_1)$. Every accepted swap strictly
reduces the global maximum UD. Since UD values are bounded below by
zero, the outer loop must terminate.

% ─────────────────────────────────────────────────────────────────────────────
\section{DoSwap --- Data Structures Updated}

When a swap of job $a$ (on $P_a$) with job $b$ (on $P_b$) is committed,
six quantities are updated atomically:

\begin{center}
\begin{tabular}{ll}
\toprule
\textbf{Structure} & \textbf{Update} \\
\midrule
$\text{mapping}[a]$ & $P_a \to P_b$ \\
$\text{mapping}[b]$ & $P_b \to P_a$ \\
$\text{proc\_jobs}[P_a]$ & remove $a$, add $b$ \\
$\text{proc\_jobs}[P_b]$ & remove $b$, add $a$ \\
$UD(P_a)$ & $-\,ud(a) + ud(b)$ \\
$UD(P_b)$ & $-\,ud(b) + ud(a)$ \\
$MU(P_a)$ & $-\,\mu(a) + \mu(b)$ \\
$MU(P_b)$ & $-\,\mu(b) + \mu(a)$ \\
\bottomrule
\end{tabular}
\end{center}

The running tallies $UD$ and $MU$ are maintained incrementally so that
Conditions~1 and~2 can be evaluated in $O(1)$ per candidate pair.
Condition~3 (DBF) is only computed after~1 and~2 pass, keeping the
typical per-pair cost low.

% ─────────────────────────────────────────────────────────────────────────────
\section{Comparison of 2a and 2b}

\begin{center}
\begin{tabular}{p{2.2cm}p{2.2cm}p{2.2cm}}
\toprule
& \textbf{Step 2a} & \textbf{Step 2b} \\
\midrule
\textbf{Structure} &
Single linear pass over all doublets &
Outer loop + inner scan, repeated \\[4pt]
\textbf{Target} &
Any two processors the current pair happens to sit on &
Always the globally highest-UD processor $P_1$ \\[4pt]
\textbf{After a swap} &
Continue to next pair (no break) &
Break inner loop; re-identify $P_1$ \\[4pt]
\textbf{Doublet scans} &
Exactly one &
One per accepted swap \\[4pt]
\textbf{UD condition} &
Pair-level imbalance $|UD(P_a)-UD(P_b)|$ decreases &
Global max UD strictly decreases \\[4pt]
\textbf{Approach} &
Distributed: fixes many local imbalances simultaneously &
Greedy: always attacks the single biggest problem first \\[4pt]
\textbf{Run time} &
$O(N^2 \cdot W^2)$ where $W$ = distinct time windows &
$O(k \cdot N^2 \cdot W^2)$, $k$ = number of accepted swaps \\
\bottomrule
\end{tabular}
\end{center}

In practice, 2a is faster but may miss the globally optimal swap order.
2b is more targeted but rescans from the beginning each time.

% ─────────────────────────────────────────────────────────────────────────────
\section{Why SPS Uses Utilisation, Not Execution Time}

SPS operates on the load value $\ell_i = e_{i,0}/p_i$, not on $e_{i,0}$
directly. Raw execution time alone is meaningless for scheduling balance:
a job with $e_m = 10$ and $p = 10$ occupies a processor 100\% of the time,
while one with $e_m = 10$ and $p = 100$ occupies only 10\%.
Utilisation correctly normalises for the period and directly measures
``fraction of processor time consumed.''

The feasibility checks in the quantum SPS mapping also work in utilisation
units: total active utilisation must not exceed remaining capacity, and the
per-processor utilisation must stay $\leq 1.0$ after each assignment.

% ─────────────────────────────────────────────────────────────────────────────
\section{Why Utility Density is the Right Metric}

Three natural alternative metrics for the refinement are compared below.

\begin{center}
\renewcommand{\arraystretch}{1.3}
\begin{tabular}{p{1.7cm}p{1.3cm}p{1.3cm}p{2.0cm}}
\toprule
\textbf{Metric} &
\textbf{Reward?} &
\textbf{Amount?} &
\textbf{Problem} \\
\midrule
$e_{i,0}/p_i$ \newline (mand.\ util.) &
No & No &
SPS already balances this --- refinement is a no-op \\
\midrule
$u_i$ only &
Yes & No &
Ignores how much optional work exists and how often the task repeats \\
\midrule
$(c_{i,N} - c_{i,0})/p_i$ \newline (opt.\ exec.) &
No & Yes &
Treats all optional work as equally valuable; misleads when $u_i$ varies \\
\midrule
$u_i\!\cdot\!(c_{i,N}-c_{i,0})/p_i$ \newline \textbf{(UD, used)} &
\textbf{Yes} & \textbf{Yes} &
Overestimates achievable; needs mutil guard \\
\bottomrule
\end{tabular}
\end{center}

UD = ``if all optional segments ran, how much utility per unit clock time
would this task contribute?'' It correctly weights tasks that are both
frequent (small $p_i$) and valuable (large $u_i$) with large optional
segments.

\textbf{Known limitation:} UD uses the theoretical maximum ($c_{i,N}$).
In practice, timing or energy may allow only a subset of optional segments.
The actual achievable utility per processor cannot be computed without
running Phase~5 first (chicken-and-egg). UD is therefore the best available
proxy computable before scheduling.

% ─────────────────────────────────────────────────────────────────────────────
\section{Full Pipeline --- Heuristic v5}

\begin{center}
\renewcommand{\arraystretch}{1.2}
\begin{tabular}{clp{4.2cm}}
\toprule
\textbf{Phase} & \textbf{Name} & \textbf{What it does} \\
\midrule
1  & Quantum SPS &
     Maps jobs to processors balancing $e_m/p_i$. Knows nothing about utility. \\
1b & \textbf{Refine Mapping} &
     \textbf{Rebalances $UD(x)$ via doublet swaps. Preserves $MU(x)$ and DBF.} \\
2  & Left Shift &
     Diagnostic: confirms timing slack at $f_{\max}$. \\
3  & Energy Slack &
     Checks if mandatory segments at $f_{\max}$ fit within budget $B$. \\
4  & Freq.\ Scaling &
     Attempts to lower frequency to save energy (no-op at $\alpha{=}1, \beta{=}0.5$). \\
5  & Greedy Optional &
     Adds optional segments greedily by utility density, checking DBF + energy. \\
6  & Pairwise Swap &
     Swaps low-value optional segments for high-value ones; fill pass after each. \\
\bottomrule
\end{tabular}
\end{center}

\begin{bluebox}
\textbf{Why Phase~1b helps Phase~5.}
Phase~5 can only schedule optional segments where there is timing slack.
Timing slack on processor $x$ is approximately proportional to
$1 - MU(x)$. If Phase~1b has spread high-$UD$ jobs across processors
while keeping $MU(x)$ balanced, then:
(a) every processor has similar timing slack available, and
(b) every processor has high-utility jobs ready to fill that slack.
Phase~5 can therefore exploit capacity on all processors, instead of
being blocked on the congested one and idle on the empty one.
\end{bluebox}

% ─────────────────────────────────────────────────────────────────────────────
\section{Concrete Worked Example}

Suppose 2 processors, 4 tasks, hyper-period $h = 10$.

\begin{center}
\small
\begin{tabular}{ccccccc}
\toprule
Task & $e_m$ & $p$ & $u$ & $e_{\text{opt}}$ & $\mu(i)$ & $ud(i)$ \\
\midrule
$T_A$ & 0.4 & 1 & 8 & 0.5 & 0.40 & 4.0 \\
$T_B$ & 0.3 & 1 & 7 & 0.4 & 0.30 & 2.8 \\
$T_C$ & 0.1 & 1 & 1 & 0.3 & 0.10 & 0.3 \\
$T_D$ & 0.1 & 1 & 1 & 0.2 & 0.10 & 0.2 \\
\bottomrule
\end{tabular}
\end{center}

\textbf{After SPS} (balances $\mu$):
\[
  P_0 = \{T_A, T_B\},\;\; MU(P_0)=0.7,\;\; UD(P_0)=6.8
\]
\[
  P_1 = \{T_C, T_D\},\;\; MU(P_1)=0.2,\;\; UD(P_1)=0.5
\]
$MU$ is imbalanced because SPS had limited choice (it also wasn't perfectly
balanced in this small example). More critically, $UD$ is catastrophically
skewed: all the utility potential is on $P_0$, which has only $30\%$ of its
time free for optional segments.

\textbf{Doublet pairs} (sorted by $|ud_a - ud_b|$, top entries):

\begin{center}
\small
\begin{tabular}{ccc}
\toprule
Pair & $|ud_a - ud_b|$ & On same proc? \\
\midrule
$(T_C, T_D)$ & 0.1 & Yes (skip) \\
$(T_A, T_B)$ & 1.2 & Yes (skip) \\
$(T_B, T_C)$ & 2.5 & No \\
$(T_B, T_D)$ & 2.6 & No \\
$(T_A, T_C)$ & 3.7 & No \\
$(T_A, T_D)$ & 3.8 & No \\
\bottomrule
\end{tabular}
\end{center}

\textbf{Step 2b} identifies $P_1 = P_0$ ($UD=6.8$). Scans doublets. First
pair with $\mathit{hi}$ on $P_0$: $(T_B, T_C)$ --- $T_B$ is on $P_0$.

Check Condition~1: swap $T_B \leftrightarrow T_C$:
\[
  UD'(P_0) = 6.8 - 2.8 + 0.3 = 4.3
\]
\[
  UD'(P_1) = 0.5 - 0.3 + 2.8 = 3.0
\]
$\max(4.3, 3.0) = 4.3 < 6.8$. \checkmark

Check Condition~2: $\mu(T_B)=0.30$, $\mu(T_C)=0.10$:
\[
  MU'(P_0)=0.7-0.3+0.1=0.5,\;\; MU'(P_1)=0.2-0.1+0.3=0.4
\]
$\delta_{\text{new}}=|0.5-0.4|=0.1 < \delta_{\text{cur}}=|0.7-0.2|=0.5$. \checkmark

Condition~3 (DBF): both processors feasible. \checkmark

\textbf{Swap committed.} New state:
\[
  P_0 = \{T_A, T_C\},\;\; MU=0.50,\;\; UD=4.3
\]
\[
  P_1 = \{T_B, T_D\},\;\; MU=0.40,\;\; UD=3.0
\]
Re-identify $P_1 = P_0$ ($UD=4.3$). Next scan finds $(T_A, T_D)$ with $T_A$
on $P_0$. Swap $T_A \leftrightarrow T_D$:
$UD'(P_0)=4.3-4.0+0.2=0.5$, $UD'(P_1)=3.0-0.2+4.0=6.8$.
$\max(0.5, 6.8)=6.8 \geq 4.3$. Condition~1 fails --- rejected. The algorithm
eventually terminates with a reasonably balanced state, having made the one
beneficial swap. Phase~5 now sees both processors with similar mandatory load
and similar utility potential.

% ─────────────────────────────────────────────────────────────────────────────
\section{Summary}

\begin{enumerate}[nosep]
  \item SPS balances mandatory utilisation $e_m/p_i$ but ignores utility.
  \item Phase~1b (Refine Mapping) fixes this by rebalancing utility density
        $ud(i) = u_i \cdot (c_{i,N} - c_{i,0}) / p_i$ across processors.
  \item Doublets (all job pairs sorted by ascending $|ud_a - ud_b|$)
        provide the candidates for swapping.
  \item Every swap must pass three conditions: UD improves, mandatory
        utilisation balance does not worsen, DBF remains feasible.
  \item \textbf{v5a} makes a single pass; \textbf{v5b} iterates,
        always targeting the most overloaded processor.
  \item The mandatory utilisation guard (Condition~2) was essential:
        without it, UD-focused swaps disrupted timing slack and reduced
        the utility that Phase~5 could achieve.
  \item UD is the best cheap proxy for ``utility potential'' available
        before scheduling, correctly combining reward rate $u_i$,
        optional volume $e_{\text{opt}}$, and task frequency $1/p_i$.
\end{enumerate}

\end{document}
